% GENERATED FILE. Edit canonical lessons, then run npm run generate:books.
% canonical-source-hash: fnv1a64:a12bbb0daf3eb85d
% canonical-lessons: PT-C492-cunhado, PT-C492-padrasto, PT-C492-madrasta, PT-C492-enteado, PT-C492-parente

\chapter{Uma família alargada}
\label{ch:pt-uma-familia-alargada}
\hlchaptermodality{492}

\begin{chapteropening}
\textbf{By the end of this chapter:} \emph{I can talk about in-laws, step-parents and relatives.}

Complete the last lesson of chapter 492: I can talk about in-laws, step-parents and relatives.
\end{chapteropening}

\section[o cunhado]{o cunhado — brother-in-law}

\label{lesson:PT-C492-cunhado}



\begin{quote}
Before the new one: say the Portuguese for to fall in love, then the Portuguese for to hug.
\end{quote}

\subsection*{You'll want to know: o cunhado}
\textbf{o cunhado} — \textquotedblleft{}brother-in-law\textquotedblright{}.

\begin{quote}
O meu cunhado é casado com a minha irmã.
\end{quote}

\begin{quote}
\emph{My brother-in-law is married to my sister.}
\end{quote}

\emph{A cunhada} — sister-in-law.

\begin{cousinweb}
\emph{cunhado} is inherited, through Old Portuguese \emph{cunnado}, from Latin \emph{cognātus} \textquotedblleft{}related by blood\textquotedblright{} (EtymDB).
\end{cousinweb}

\begin{grammarlens}[title={What you've built}]
One more word for telling, explaining and describing at B1.
\end{grammarlens}

\subsection*{Your turn}
\begin{itemize}
\raggedright
  \item \emph{Say it:} \emph{o cunhado}
  \item \emph{Say it:} \emph{o cunhado}, once more
  \item \emph{Say it:} a sentence of your own with \emph{o cunhado}
  \item \emph{Recall:} say the Portuguese for to fall in love, then the Portuguese for to hug, then say \emph{o cunhado} again
\end{itemize}

\begin{tcolorbox}[breakable,colback=teal!4,colframe=teal!35!black,title={Before you move on}]
What does o cunhado mean? (\textquotedblleft{}Brother-in-law\textquotedblright{}.) Say it once more.
\end{tcolorbox}
\section[o padrasto]{o padrasto — stepfather}

\label{lesson:PT-C492-padrasto}



\begin{quote}
Before the new one: say the Portuguese for to hug, then the Portuguese for brother-in-law.
\end{quote}

\subsection*{You'll want to know: o padrasto}
\textbf{o padrasto} — \textquotedblleft{}stepfather\textquotedblright{}.

\begin{quote}
Ela dá-se muito bem com o padrasto.
\end{quote}

\begin{quote}
\emph{She gets on very well with her stepfather.}
\end{quote}

Built on Latin \emph{pater} \textquotedblleft{}father\textquotedblright{}, not on \emph{pai}.

\begin{cousinweb}
\emph{padrasto} is inherited from Latin \emph{patraster} (EtymDB), built on \emph{pater} \textquotedblleft{}father\textquotedblright{}; in Latin inscriptions the word named a father-in-law (Lewis \& Short), and in Portuguese it settled on a stepfather.
\end{cousinweb}

\begin{grammarlens}[title={What you've built}]
One more word for telling, explaining and describing at B1.
\end{grammarlens}

\subsection*{Your turn}
\begin{itemize}
\raggedright
  \item \emph{Say it:} \emph{o padrasto}
  \item \emph{Say it:} \emph{o padrasto}, once more
  \item \emph{Say it:} a sentence of your own with \emph{o padrasto}
  \item \emph{Recall:} say the Portuguese for to hug, then the Portuguese for brother-in-law, then say \emph{o padrasto} again
\end{itemize}

\begin{tcolorbox}[breakable,colback=teal!4,colframe=teal!35!black,title={Before you move on}]
What does o padrasto mean? (\textquotedblleft{}Stepfather\textquotedblright{}.) Say it once more.
\end{tcolorbox}
\section[a madrasta]{a madrasta — stepmother}

\label{lesson:PT-C492-madrasta}



\begin{quote}
Before the new one: say the Portuguese for brother-in-law, then the Portuguese for stepfather.
\end{quote}

\subsection*{You'll want to know: a madrasta}
\textbf{a madrasta} — \textquotedblleft{}stepmother\textquotedblright{}.

\begin{quote}
A minha madrasta trata-me como uma filha.
\end{quote}

\begin{quote}
\emph{My stepmother treats me like a daughter.}
\end{quote}

A neutral word in everyday use, though fairy tales made \emph{a madrasta má} a stock villain.

\begin{cousinweb}
\emph{madrasta} is usually traced to Late Latin \emph{mātrastra} \textquotedblleft{}stepmother\textquotedblright{}, built on \emph{māter} \textquotedblleft{}mother\textquotedblright{} (Lewis \& Short).
\end{cousinweb}

\begin{grammarlens}[title={What you've built}]
One more word for telling, explaining and describing at B1.
\end{grammarlens}

\subsection*{Your turn}
\begin{itemize}
\raggedright
  \item \emph{Say it:} \emph{a madrasta}
  \item \emph{Say it:} \emph{a madrasta}, once more
  \item \emph{Say it:} a sentence of your own with \emph{a madrasta}
  \item \emph{Recall:} say the Portuguese for brother-in-law, then the Portuguese for stepfather, then say \emph{a madrasta} again
\end{itemize}

\begin{tcolorbox}[breakable,colback=teal!4,colframe=teal!35!black,title={Before you move on}]
What does a madrasta mean? (\textquotedblleft{}Stepmother\textquotedblright{}.) Say it once more.
\end{tcolorbox}
\section[o enteado]{o enteado — stepson}

\label{lesson:PT-C492-enteado}



\begin{quote}
Before the new one: say the Portuguese for stepfather, then the Portuguese for stepmother.
\end{quote}

\subsection*{You'll want to know: o enteado}
\textbf{o enteado} — \textquotedblleft{}stepson\textquotedblright{}.

\begin{quote}
O meu enteado tem doze anos.
\end{quote}

\begin{quote}
\emph{My stepson is twelve.}
\end{quote}

\emph{A enteada} — stepdaughter.

\begin{cousinweb}
\emph{enteado} goes back to Old Portuguese \emph{entẽado} (EtymDB) and is usually traced to Latin \emph{ante nātum} \textquotedblleft{}born before\textquotedblright{}: a child from an earlier marriage.
\end{cousinweb}

\begin{grammarlens}[title={What you've built}]
One more word for telling, explaining and describing at B1.
\end{grammarlens}

\subsection*{Your turn}
\begin{itemize}
\raggedright
  \item \emph{Say it:} \emph{o enteado}
  \item \emph{Say it:} \emph{o enteado}, once more
  \item \emph{Say it:} a sentence of your own with \emph{o enteado}
  \item \emph{Recall:} say the Portuguese for stepfather, then the Portuguese for stepmother, then say \emph{o enteado} again
\end{itemize}

\begin{tcolorbox}[breakable,colback=teal!4,colframe=teal!35!black,title={Before you move on}]
What does o enteado mean? (\textquotedblleft{}Stepson\textquotedblright{}.) Say it once more.
\end{tcolorbox}
\section[o parente]{o parente — relative}

\label{lesson:PT-C492-parente}



\begin{quote}
Before the new one: say the Portuguese for stepmother, then the Portuguese for stepson.
\end{quote}

\subsection*{You'll want to know: o parente}
\textbf{o parente} — \textquotedblleft{}relative\textquotedblright{}.

\begin{quote}
Nas festas de família juntamos todos os parentes.
\end{quote}

\begin{quote}
\emph{At family parties we get all our relatives together.}
\end{quote}

Not English \emph{parent}: it means any relative. Parents are \emph{os pais}.

\begin{cousinweb}
\emph{parente} is inherited from Latin \emph{parentem} \textquotedblleft{}a parent\textquotedblright{} (EtymDB), from \emph{parere} \textquotedblleft{}to give birth\textquotedblright{}; in Portuguese it widened to any relative.
\end{cousinweb}

\begin{grammarlens}[title={What you've built}]
One more word for telling, explaining and describing at B1.
\end{grammarlens}

\subsection*{Your turn}
\begin{itemize}
\raggedright
  \item \emph{Say it:} \emph{o parente}
  \item \emph{Say it:} \emph{o parente}, once more
  \item \emph{Say it:} a sentence of your own with \emph{o parente}
  \item \emph{Recall:} say the Portuguese for stepmother, then the Portuguese for stepson, then say \emph{o parente} again
\end{itemize}

\begin{tcolorbox}[breakable,colback=teal!4,colframe=teal!35!black,title={Before you move on}]
What does o parente mean? (\textquotedblleft{}Relative\textquotedblright{}.) Say it once more.
\end{tcolorbox}
